\documentclass[10pt,a4paper]{amsart}
\usepackage[margin=19mm]{geometry}
\usepackage{amsmath,amssymb,mathtools}
\usepackage[scr=rsfs,cal=euler]{mathalfa}
\usepackage{xfrac}
\usepackage[unicode,hidelinks,bookmarks=false]{hyperref}
\newcommand{\ie}{i.e.\ }
\newcommand{\cf}{cf.\ }
\newcommand{\resp}{respectively\ }
\newcommand{\Subsection}{\subsection}
\providecommand{\nobreakdash}{\nobreak\hbox{-}}
\setlength{\emergencystretch}{2em}
\multlinegap0pt
\usepackage{tikz}
\newtheorem{lemma}{Lemma}
\newtheorem{theorem}{Theorem}
\title{On Threshold Pairwise Compatibility Graphs}
\author{Sheikh Azizul Hakim, Md. Shamsuzzoha Bayzid}
\date{Source: arXiv:2604.20042v2 (CC BY 4.0)}
\begin{document}
\maketitle

\noindent\emph{Source and licence.} On Threshold Pairwise Compatibility Graphs. Version arXiv:2604.20042v2 (CC BY 4.0), distributed by arXiv under the Creative Commons Attribution 4.0 International licence, http://creativecommons.org/licenses/by/4.0/, as recorded in the arXiv licence metadata for this exact version and shown on the abstract page https://arxiv.org/abs/2604.20042v2. Full licence notice: \texttt{licenses/arxiv-2604.20042-CC-BY-4.0.txt}.\par\smallskip

\noindent\textit{Editorial scope.} Counted unit: the article's theorem thm:threshold-set-system (source0.tex lines 1156--1164). The article's complete original proof of it is delivered with it (source0.tex lines 1166--1246, one \texttt{proof} environment of 81 lines). No mathematics is written, repaired or re-derived: the statement and the proof are the article's own lines, in the article's own notation, and no retained line is substituted or re-flowed.  Retained uncounted as the source-local context the proof needs: lines 1084--1100.  Nothing was omitted from any retained span, so the package carries no omission record.  No retained line contains a citation, so the wrapper carries no bibliography and no bibliography entry.  No retained line contains a figure environment, an image-inclusion command or drawing code, so no image is delivered and the package declares no asset dependency.  Editorial formatting only, in addition: an AMS \texttt{amsart} preamble that restates the article's own declarations and macros used by the retained lines, namely the environments \texttt{lemma}, \texttt{theorem} on separate counters, together with the standard packages those lines need.  The retained environments print in this document as Lemma 1, Theorem 1 in order of appearance, whereas the article prints the same items with the numbering of its own full document.  The complete native source of the article as distributed in the arXiv e-print for this exact version is bundled unchanged as \texttt{sources/198899/source0.tex}.
\begin{lemma}[Single-shell bound]\label{lem:single-shell-bound}
There exists an absolute constant \(C>0\) such that the following holds.

Let \(T\) be a reduced weighted tree with leaf set \(A\), where \(|A|=m\), and let
\(I=[\alpha,\beta]\) be an interval of real numbers. Define
\[
\mathcal F(T,I):=
\Bigl\{
\{a\in A:d_T(x,a)+\lambda\in I\}: x\in T,\ \lambda\in\mathbb R
\Bigr\},
\]
where \(x\in T\) means that \(x\) is an arbitrary point of the geometric realization of
the tree \(T\). Then
\(
|\mathcal F(T,I)|\le C\,m^3.
\)
\end{lemma}

\begin{theorem}\label{thm:threshold-set-system}
There exists an absolute constant $C>0$ such that the following holds.

Let $A$ be a finite set with $|A|=m$, and let $\mathcal S\subseteq 2^A$. If
$G(A,\mathcal S)$ is a \textsc{$(k,t)$-threshold-\textsc{PCG}}, then
\(
|\mathcal S|\le (C\,m^3)^k.
\)
\end{theorem}

\begin{proof}
Assume that $G(A,\mathcal S)$ is a \textsc{$(k,t)$-threshold-\textsc{PCG}}. Thus there
exist \textsc{PCG}s
\(
G_1,\dots,G_k
\)
on the common vertex set $A\cup B$ such that, for every $a\in A$ and every
$S\in\mathcal S$,
\[
ab_S\in E(G(A,\mathcal S))
\iff
\sum_{i=1}^k \mathbf 1_{\{ab_S\in E(G_i)\}} \ge t.
\]

Fix $i\in[k]$. Since $G_i$ is a \textsc{PCG}, there exist a weighted tree
$\widetilde T_i$, an interval $I_i=[\alpha_i,\beta_i]$, and a leaf-labeling bijection
such that
\(
xy\in E(G_i)
\iff
d_{\widetilde T_i}(x,y)\in I_i
\)
for all distinct leaves corresponding to vertices of $A\cup B$.

Let $T_i$ be the minimal subtree of $\widetilde T_i$ spanning the leaves corresponding
to $A$, and suppress all degree-$2$ vertices. Then $T_i$ is a reduced weighted tree with
leaf set exactly $A$.

For each $S\in\mathcal S$, let $x_{i,S}$ be the unique closest point of $T_i$ to the
leaf corresponding to $b_S$, and let
\(
\lambda_{i,S}:=d_{\widetilde T_i}(b_S,x_{i,S}).
\)
Then for every $a\in A$,
\(
d_{\widetilde T_i}(b_S,a)=\lambda_{i,S}+d_{T_i}(x_{i,S},a),
\)
and hence
\(
ab_S\in E(G_i)
\iff
d_{T_i}(x_{i,S},a)+\lambda_{i,S}\in I_i.
\)

Define
\(
\mathcal F_i:=
\Bigl\{
\{a\in A:d_{T_i}(x,a)+\lambda\in I_i\}: x\in T_i,\ \lambda\in\mathbb R
\Bigr\}.
\)
By Lemma~\ref{lem:single-shell-bound},
\(
|\mathcal F_i|\le C\,m^3
\ \text{for each }i\in[k].
\)

Now fix $S\in\mathcal S$, and let
\(
X_{i,S}:=\{a\in A:ab_S\in E(G_i)\}.
\)
Then $X_{i,S}\in\mathcal F_i$ for each $i\in[k]$, and by the threshold rule,
\(
S=
\Bigl\{
a\in A:
\sum_{i=1}^k \mathbf 1_{\{a\in X_{i,S}\}} \ge t
\Bigr\}.
\)
Therefore, each set $S\in\mathcal S$ is determined by an $k$-tuple
\(
(X_{1,S},\dots,X_{k,S})\in \mathcal F_1\times\cdots\times\mathcal F_k.
\)
Hence,
\[
|\mathcal S|
\le
|\mathcal F_1|\cdots |\mathcal F_k|
\le
(C\,m^3)^k.
\] \end{proof}

\end{document}
